% 制作封面
% 加载宏包的命令在format之中，这里为了完整，单独使用时需要利用这里的有关命令
% \usepackage[svgnames]{xcolor} % Required to specify font color
% \newcommand*{\plogo}{\fbox{$\mathcal{PL}$}} % Generic publisher logo 
% \usepackage{graphicx} % Required for box manipulation

%------------------------------------------------------------------------------
%	TITLE PAGE
%------------------------------------------------------------------------------


\newcommand*{\rotrt}[1]{\rotatebox{90}{#1}} % Command to rotate right 90 degrees
\newcommand*{\rotlft}[1]{\rotatebox{-90}{#1}} % Command to rotate left 90 degrees

\newcommand*{\titleBC}{\begingroup % Create the command for including the title page in the document
\centering % Center all text
\def\CP{\textit{\Huge 代数$K$理论讨论班笔记}} % Title
\settowidth{\unitlength}{\CP} % Set the width of the curly brackets to the width of the title
{\color{LightGoldenrod}\resizebox*{\unitlength}{\baselineskip}{\rotrt{$\}$}}} \\[\baselineskip] % Print top curly bracket
\textcolor{Sienna}{\CP} \\[\baselineskip] % Print title

{\color{RosyBrown}\Large 中国科学院大学~~数学科学学院} \\ % Tagline or further description
{\color{RosyBrown}\Large \kaishu{H. Zhang}}\\
{\color{LightGoldenrod}\resizebox*{\unitlength}{\baselineskip}{\rotlft{$\}$}}} % Print bottom curly bracket
\vfill % Whitespace between the title and the author name

\begin{center}
\frame{\includegraphics[width=90mm]{./figures/seminar.jpg}}\\ %载入封面图片

{\large 从左至右依次为\\
Quillen~~Milnor~~Grothendieck}\\
\vfill
% \plogo\\[0.5\baselineskip] %plogo未使用
\today
\end{center}
\endgroup}